
\subsubsection*{4.4.2}
\phantomsection
\label{sga3:VIIB:4.4.2}

Continue to assume that \(k\) is a field of characteristic \(p\).
Let \(G\) be an infinitesimal formal \(k\)-group, with affine
algebra \(A\). Since \(G\) is infinitesimal, every open ideal of
\(A\) contains
\[
  I_A^{\{p^n\}}
\]
for all sufficiently large \(n\). Thus \(A\) is the projective
limit of the affine algebras
\[
  A/I_A^{\{p^n\}}
\]
of the groups \(\operatorname{Fr}_nG\)
\textup{(cf.\ \SGARefLink{sga3:VIIA:4.1.2}{4.1.2})}. Equivalently, every infinitesimal formal
\(k\)-group is an inductive limit of formal \(k\)-groups of finite
height.

Suppose that \(G\) has height at most \(n\), and write
\[
  H=G/\operatorname{Fr}G.
\]
\footnote{\label{sga3:VIIB:nde:141}The editors supplied the details
of the original argument in what follows.}
By \SGARefLink{sga3:VIIB:2.4}{2.4} and \SGARefLink{sga3:VIIB:2.4.1}{2.4.1}, the morphism
\[
  \operatorname{Fr}:G\longrightarrow G^{(p)}
\]
factors as an epimorphism
\[
  \pi:G\longrightarrow H
\]
followed by a monomorphism
\[
  i:H\longrightarrow G^{(p)}.
\]
We therefore have the following commutative diagram:
\begin{center}
  \phantomsection
  \label{sga3:VIIB:4.4.2:height-reduction-frobenius-diagram}
  \begin{tikzcd}[row sep=2.8em,column sep=3em]
G \arrow[r,"\pi"] \arrow[dr,"\operatorname{Fr}"']
  & H \arrow[r,"\operatorname{Fr}^{\,n-1}"] \arrow[d,"i"]
  & H^{(p^{\,n-1})}
      \arrow[d,"i^{(p^{\,n-1})}" right] \\
{} & G^{(p)} \arrow[r,"\operatorname{Fr}^{\,n-1}"]
  & G^{(p^n)}
\end{tikzcd}

\end{center}
Since the functor \(X\mapsto X^{(p)}\) commutes with fiber
products, the morphism
\[
  i^{(p^{\,n-1})}
\]
is again a monomorphism. Moreover,
\(\operatorname{Fr}^{n}(G/k)\) is zero and \(\pi\) is an
epimorphism. Hence
\[
  \operatorname{Fr}^{n-1}\bigl(G^{(p)}/k\bigr)\circ i
  =
  i^{(p^{\,n-1})}\circ
  \operatorname{Fr}^{n-1}(H/k)
\]
is zero. Because \(i^{(p^{\,n-1})}\) is a monomorphism,
\[
  \operatorname{Fr}^{n-1}(H/k)=0,
\]
so \(H\) has height at most \(n-1\).

It follows that every formal \(k\)-group of finite height has a
composition series whose quotients have height at most \(1\), and
therefore can be described by \(p\)-Lie algebras over \(k\).

Finally, the affine algebra \(A\) of \(G\) is a quotient of
\[
  k[[\omega_{G/k}]]
\]
by \SGARefLink{sga3:VIIB:4.4.A}{Remark~4.4.A}. If \(\omega_{G/k}\) is finite-dimensional over
\(k\), all the algebras
\[
  A/I_A^{\{p^n\}}
\]
are finite-dimensional \(k\)-vector spaces. Consequently, every
infinitesimal formal group \(G\) over a field of characteristic
\(p>0\) for which \(\omega_{G/k}\) is finite-dimensional over
\(k\) is an inductive limit of finite formal groups
\textup{(that is, groups of finite length; cf.\ \SGARefLink{sga3:VIIB:1.2.6}{1.2.6})}.
